\chapter*{Abstract}
\addcontentsline{toc}{chapter}{Abstract}

This Assignment 2 implementation delivers a working collaborative document editor with a React frontend,
FastAPI backend, authenticated realtime collaboration, and a streaming AI writing assistant. The final
system supports registration, login, refresh-token based session recovery, local password reset,
document CRUD, role-based sharing, version history and restore, autosave, live collaboration, AI rewrite,
summarize, and translate operations, AI history, share-by-link, and PDF/DOCX export.

The implementation is intentionally organised as a monorepo with clear service boundaries. The FastAPI
backend owns authentication, permissions, document lifecycle, AI orchestration, export workflows, and
realtime session issuance. The React client provides the user-facing editing, sharing, AI review, and
history flows. A dedicated Yjs/WebSocket relay handles low-latency collaborative state propagation. This
division keeps the codebase understandable, testable, and aligned with the architecture described in
Assignment 1 while still allowing pragmatic implementation changes.

Several design decisions were driven by local reliability and reproducibility. Passwords are securely
hashed, refresh tokens are hashed before persistence, and websocket authentication uses signed
subprotocol tokens rather than query-string credentials. AI suggestions are streamed progressively,
support cancellation, and remain subject to explicit user review before application. Under concurrent
editing, stale AI suggestions are blocked from overwriting text that has already changed. These choices
prioritise both security and collaboration correctness rather than a narrow feature checklist.

The final repository also includes meaningful automated verification: backend tests, frontend component
tests, realtime conflict-resolution tests, a root integration smoke test, and browser-level Playwright
coverage. Alongside the implementation, the repository contains a README, a deviation report, API
contracts, and setup scripts so that the system can be cloned, configured, and run locally without
undocumented steps.

\clearpage
